\section{Introduction}
Understanding transition has been a long-standing problem in fluid mechanics since the pipe-flow experiments of \citet{reynolds-1883} and has continued to vex generations of scientists seeking to unravel the mechanisms leading to turbulence. A cornerstone of this effort has been linear stability theory, in which the governing equations are linearized about a fixed base state and the system is classified as linearly stable or unstable through the eigenspectrum of the resulting operator. This approach has provided the foundation for the study of hydrodynamic instabilities, including the Orr--Sommerfeld problem for parallel shear flows~\citep{orr-1907,sommerfeld-1908} and Taylor--Couette flow between rotating cylinders~\citep{taylor-1923}. These and many other examples are collected in the classical texts of \citet{chandrasekhar-1961} and \citet{drazin-reid-2004}.

Despite its broad success, and its continued extension to complex two- and three-dimensional flows~\citep{theofilis-2011}, the predictive limitations of modal stability analysis became apparent early on, particularly for parallel shear flows. In fact, for plane Couette flow, the system is predicted to remain linearly stable for all Reynolds numbers~\citep{romanov-1973}, even though transition is observed in experiments.


To address this discrepancy, subsequent developments in hydrodynamic stability theory shifted the focus beyond classical eigenvalue analysis toward non-normal stability theory, where the non-orthogonality of the eigenvectors was recognized as a mechanism capable of producing substantial transient amplification, even when all eigenmodes are asymptotically stable~\citep{trefethen-1993}. Disturbances may therefore grow sufficiently to trigger nonlinear effects and transition well before the late-time dynamics predicted by the eigenspectrum become relevant. By identifying the optimal energy amplification attainable over all possible initial perturbations, non-modal stability theory has found considerable success in describing transition in boundary layers~\citep{hoepffner-brandt-henningson-2005,monokrousos-et-al-2010} and pipe flows~\citep{meseguer-trefethen-2003}, and many other set-ups. A comprehensive review of the theoretical framework and applications is provided by \citet{schmid-2007}.

Although many of these analyses considered perturbations evolving about steady base states, the initial-value formulation underlying non-modal stability theory is not restricted to autonomous systems. It extends naturally to flows whose base state evolves in time~\citep{schmid-2007}, enabling transient amplification to be investigated in a broad range of time-dependent configurations. Examples include temporally evolving mixing layers undergoing Kelvin--Helmholtz roll-up and pairing~\citep{arratia-caulfield-chomaz-2013,kaminski-caulfield-taylor-2014}, oscillatory and solitary-wave boundary layers~\citep{biau-2016,verschaeve-pedersen-tropea-2018}, and pulsatile internal flows~\citep{xu-song-avila-2021,pier-schmid-2021}. Nevertheless, optimal-growth analysis remains intrinsically tied to a prescribed initial time and optimization horizon. It therefore identifies the disturbance that maximizes amplification over a particular interval, rather than providing a continuously evolving description of the dominant stability mechanisms along the underlying base-state trajectory.

A framework that addresses this limitation was introduced by \citet{babaee-sapsis-2016} through optimally time-dependent (OTD) modes. The approach employs a minimization principle to construct a finite-dimensional, time-dependent orthonormal basis that evolves along the underlying trajectory and spans the dynamically most unstable directions of the tangent space.

\citet{kern-et-al-2021} applied this method to pulsating Poiseuille flow and identified subharmonic perturbation cycles that would remain hidden if the disturbances were restricted to the periodicity of the base flow, as in Floquet theory~\citep{davis-1976}. A related application was presented by \citet{beneitez-et-al-2023}, who investigated the finite-time stability of a fully three-dimensional, unsteady trajectory in the Blasius boundary layer initiated by the minimal seed~\citep{pringle-kerswell-2010,cherubini-et-al-2011}. By using OTD modes to construct a reduced linearized operator along this trajectory, they showed that meaningful instantaneous stability information such as finite-time Lyapunov exponents and coherent flow structures could be extracted even from a strongly aperiodic base flow. More recently, this strategy has been used to examine the onset of global instability in configurations of direct relevance to external aerodynamics such as the flow over a pitching airfoil~\citep{kern-et-al-2024}.

Despite these promising applications, the OTD framework has several acknowledged limitations~\citep{kern-et-al-2021}. The dimension \(r\) of the OTD subspace must be prescribed \emph{a priori}, and no general criterion is available for determining how many modes should be retained. Excluding relevant stable directions may cause the reduced system to underestimate both non-normality and energy growth, particularly in strongly non-normal flows. Increasing the subspace dimension can alleviate this problem, but also increases the computational cost. Because the method requires, in principle, the solution of a full three-dimensional problem for each mode, its implementation may involve \(r+1\) coupled direct numerical simulations and become prohibitively expensive as the number of retained modes increases~\citep{beneitez-et-al-2023}. A further limitation concerns the initialization of the OTD basis, for which no general procedure is currently available~\citep{kern-et-al-2024}. Consequently, the short-time dynamics may remain strongly dependent on the chosen initialization until the basis aligns with the dynamically dominant subspace.

In summary, classical non-modal growth theory does not provide a continuously evolving description along the base-state trajectory, but only the maximal gain attainable over a prescribed time horizon. OTD modes overcome this restriction, yet introduce computational compromises that may neglect relevant non-normal behaviour and remain strongly dependent on how the basis is initialized, a choice that is itself tied to the initial conditions imposed on the system.

Another fundamental shortcoming of non-normal optimal-growth theory emerges in the context of diffusive instabilities, particularly the Rayleigh--Taylor instability (RTI), which occurs when two superposed fluids of distinct densities are subjected to an acceleration directed against their density stratification~\citep{rayleigh-1883,taylor-1950}, and will be the main focus of this paper. Among the several configurations in which this instability arises, the diffusive RTI encountered during carbon dioxide (\(\mathrm{CO}_2\)) sequestration in deep geological formations~\citep{orr-2009} provides an especially clear example of the limitations of non-modal growth theory. In this physical configuration, dissolved \(\mathrm{CO}_2\) diffuses into the underlying brine, producing a dense diffusive boundary layer above a lighter fluid. Although the resulting stratification is gravitationally unstable, the earliest dynamics remain dominated by molecular diffusion, since the convective motion has not yet developed, and the base state continues to broaden as diffusion proceeds.

\citet{riaz-et-al-2006} analyzed this onset problem using the quasi-steady-state approximation (QSSA), which is commonly employed for this class of problems. In the QSSA, the diffusive base state is frozen at a prescribed time, and a classical eigenvalue analysis is performed with time treated as a parameter. The approximation assumes that the base state evolves much more slowly than the instability and may therefore be regarded as fixed during perturbation growth~\citep{tan-homsy-1986}. This separation of time scales is, however, violated during the early diffusion-dominated stage. Later studies therefore sought to overcome the QSSA, and in particular \citet{rapaka-et-al-2008} used non-modal stability theory to compute the maximum finite-time amplification over all admissible initial perturbations. Despite this attempt, \citet{don-nils-amir-2013} showed that convection is instead more likely to be triggered by suboptimal perturbations than by the mathematically optimal disturbance obtained from the unconstrained amplification problem.

This suggests that mathematical optimality does not necessarily correspond to a physically realizable initial condition.

Meaningful predictions must instead consider the complete initial-value problem initialized with physically relevant disturbances, requiring an assessment of the effect of initial conditions. An instructive example is provided by \citet{daniel-riaz-tchelepi-2015}, who examined how a permeability field varying periodically across the aquifer thickness affects the onset of natural convection. They used ensembles of random initial perturbations localized within the diffusive boundary layer to characterize the resulting distribution of instantaneous growth rates. The ensemble-mean growth agreed closely with that obtained from the boundary-layer-localized perturbation profile of \citet{riaz-et-al-2006}. Nevertheless, there is no guarantee that this agreement would persist under a different correlation structure of the initial perturbations or for fully three-dimensional initial conditions. Consequently, the statistical specification of the initial disturbances must itself be regarded as an essential part of the stability problem.

A related, but distinct, manifestation of this sensitivity was identified by \citet{prathama-pantano-2021} for a baroclinically generated, variable-density shear layer with an analogous diffusively evolving base state \citep{gat-et-al-2017}. Even within a fully non-modal treatment, the early disturbance evolution remained dominated by the diffusive spreading of the base state, with genuine convective growth emerging only at later times. Comparison with direct numerical simulations further showed that the wavenumber assumed to be representative of the disturbance did not drive the observed growth. Instead, the growth was triggered by a subharmonic wavenumber already present in the broadband initial condition.


The available methods for analyzing transient instabilities therefore offer different advantages, but each leaves an important part of the problem unresolved. For transient instabilities, classical modal theory takes the form of the QSSA, and cannot represent the simultaneous evolution of the base state and the perturbations. Non-modal optimal-growth theory retains this evolution and captures transient amplification, but remains tied to a prescribed initial time and optimization horizon, and the resulting optimal perturbation need not be physically realizable or representative of what is actually observed, as demonstrated for both porous-media RTI and a diffusively evolving shear layer. OTD modes overcome the fixed-horizon restriction by continuously tracking the dynamically dominant directions along the trajectory. They nevertheless require a prescribed subspace dimension, depend strongly on the initialization of the basis, and incur a computational cost that increases with the number of retained modes.

It is within this context that we position the present work. We seek a description that retains the non-autonomous evolution of disturbances along a time-dependent base state while explicitly accounting for the statistics of physically admissible initial conditions, rather than isolating a single optimally amplified disturbance. Such a description should provide not only an integrated measure of amplification, but also the spatially resolved statistics needed to explain how the instability mechanism develops as the base state evolves.


The starting point is the work of \citet{frame-towne-2024}, who showed that the large optimal gains commonly invoked to explain transient amplification in non-normal systems correspond to rare events when disturbances are instead drawn from a random ensemble. They further showed that the resulting mean amplification depends on the spatial correlation of the imposed initial conditions. This result motivates a shift from the deterministic question of which individual disturbance experiences the greatest gain to the statistical question of what amplification is expected when the system is perturbed by an ensemble of initial conditions. The literature discussed above suggests that this is a natural perspective from which to analyze transient instabilities. Nevertheless, the original formulation was developed for a fixed base state and left unresolved the physical specification of the initial correlations.
